\documentclass[letterpaper,11pt]{article}

\usepackage[utf8]{inputenc}
\usepackage[T1]{fontenc}
\usepackage{lmodern}
\usepackage[top=0.35in, bottom=0.35in, left=0.55in, right=0.55in]{geometry}
\usepackage{titlesec}
\usepackage{enumitem}
\usepackage{hyperref}

% ---- Hide all link colors ----
\hypersetup{hidelinks}

% ---- Page Setup ----
\pagestyle{empty}
\setlength{\parindent}{0pt}
\setlength{\parskip}{0pt}

% ---- Custom section command for precise spacing ----
\newcommand{\ressection}[1]{%
  \vspace{3pt}%
  \noindent{\large\bfseries\scshape #1}%
  \par\vspace{1pt}%
  \hrule height 0.4pt%
  \vspace{2pt}%
}

\begin{document}

% ====================
% HEADER
% ====================
\begin{center}
  {\LARGE \textbf{Jiefu (Jeff) Zhang}} \\[2pt]
  \href{mailto:zhan4018@purdue.edu}{zhan4018@purdue.edu} $|$
  \href{https://jefferyz001.github.io}{jefferyz001.github.io} $|$
  \href{https://github.com/jefferyz001}{github.com/jefferyz001} $|$
  West Lafayette, IN \\[2pt]
  \textit{Research focus: AI for chip design and manufacturing automation (EDA)}
\end{center}

\vspace{-6pt}

% ====================
% EDUCATION
% ====================
\ressection{Education}

\begin{tabular*}{\textwidth}{@{}l@{\extracolsep{\fill}}r@{}}
  \textbf{Purdue University} & West Lafayette, IN \\
  \textit{M.S. Computer Engineering (Thesis Track)} & \textit{Expected Dec 2026} \\
  \textit{B.S. Computer Engineering} & \textit{Dec 2024} \\
\end{tabular*}

\begin{itemize}[leftmargin=1.2em, topsep=2pt, itemsep=1pt, parsep=0pt]
  \item Accelerated BS/MS Program \,|\, Graduate GPA: 3.88/4.0
  \item Coursework: ASIC Design, VLSI Design, Semiconductor Fabrication, RISC-V Architecture, Embedded Systems, Artificial Intelligence
\end{itemize}

% ====================
% RESEARCH EXPERIENCE
% ====================
\ressection{Research Experience}

\begin{tabular*}{\textwidth}{@{}l@{\extracolsep{\fill}}r@{}}
  \textbf{Graduate Research Assistant, CLAN Lab, Purdue University} & Oct 2024 -- Present \\
  \textbf{MS Student Member, Wan Lab, Columbia University} & Jun 2026 -- Present \\
\end{tabular*}
\textit{Advisors: Prof.\ Vaneet Aggarwal (Purdue) and Prof.\ Zishen Wan (Columbia)}

\vspace{2pt}
\begin{itemize}[leftmargin=1.2em, topsep=3pt, itemsep=3pt, parsep=0pt]
  \item \textbf{CARVE: Chip Placement with Reusable Expertise} (first author, \href{https://arxiv.org/abs/2609.33106}{arXiv:2609.33106}): Adapted pretrained RL macro-placement policies to new circuits without the original training data, reusing saved specialists only after they pass local validation. A reuse-first variant cut recorded training time by \textbf{58.5\%} (9.66 $\rightarrow$ 4.01 h) while mean HPWL gain moved only from 8.41\% to 7.86\%; imported specialists were reused on \textbf{6 of 7} new IBM circuits.
  \item \textbf{SIREN: Correcting the Winner's Curse in LLM Evaluation} (first author, \textbf{NeurIPS 2026}, \href{https://arxiv.org/abs/2605.05973}{arXiv}): Co-developed a selection-aware repeated-split reporting protocol with bootstrap confidence bands; led MMLU-Pro tuning experiments with vLLM-served open-weight models and DSPy. Now extending it to AI-for-EDA results (macro placement).
  \item \textbf{Safe Social Robot Navigation} (first author, \href{https://arxiv.org/abs/2603.06729}{arXiv:2603.06729}, \href{https://github.com/jefferyz001/PSS-Social}{code}): Built a PPO-based navigation framework with density-adaptive proxemic reward shaping; achieved \textbf{86\%+ safe success at 21-agent crowd densities} where DS-RNN and SARL baselines collapse, benchmarked against ORCA and SFM.
  \item \textbf{Continual and Decentralized Learning} (\href{https://arxiv.org/abs/2609.33725}{SPHERE}, \href{https://arxiv.org/abs/2605.08681}{Core-Halo}): Co-developed representation-guided replay allocation for online continual learning; built SUMO/PettingZoo multi-intersection traffic-control benchmarks and a decentralized DQN pipeline for decentralized fixed-point learning.
\end{itemize}

% ====================
% SELECTED PROJECTS
% ====================
\ressection{Selected Hardware Projects}

\begin{itemize}[leftmargin=1.2em, topsep=2pt, itemsep=3pt, parsep=0pt]
  \item \textbf{Multicore RISC-V Processor:} Designed a 5-stage pipelined RISC-V processor with hazard resolution via forwarding/stalling; improved multicore efficiency by \textbf{30\%} using MOESI cache coherence.
  \item \textbf{USB 1.1 ASIC Data Transfer:} Achieved \textbf{12 Mbps} throughput (USB 1.1 standard) in Verilog with AHB-Lite interface; validated via QuestaSim with \textbf{98\% test coverage}.
  \item \textbf{8-bit Wallace Tree Multiplier} (ECE 559 MOS VLSI Design): Full-custom schematic/layout, DRC/LVS clean; worst-case delay \textbf{$\leq$ 2.5\,ns}, energy $\leq$ 850\,fJ with parasitic extraction.
\end{itemize}

% ====================
% TECHNICAL SKILLS
% ====================
\ressection{Technical Skills}

\begin{itemize}[leftmargin=1.2em, topsep=2pt, itemsep=2pt, parsep=0pt]
  \item \textbf{Programming:} Python, C/C++, Verilog, SystemVerilog, Java, Bash, Git, \LaTeX
  \item \textbf{ML / AI for EDA:} PyTorch, RL (PPO, DQN), RL-based macro placement, continual learning, LLM evaluation, DSPy, vLLM, PettingZoo, Multi-GPU Training
  \item \textbf{Hardware/EDA:} Cadence Virtuoso, QuestaSim, DRC/LVS, parasitic extraction, RTL Design, FPGA, ASIC/SoC Integration, Linux
\end{itemize}

% ====================
% ADDITIONAL
% ====================
\ressection{Additional}

\begin{itemize}[leftmargin=1.2em, topsep=2pt, itemsep=2pt, parsep=0pt]
  \item \textbf{Award:} NeurIPS 2026 Financial Assistance Award (full conference and workshop registration and hotel stay)
  \item \textbf{Teaching:} Graduate TA, Purdue ECE 33700 -- ASIC Design Laboratory, 60+ students (Spring 2025)
  \item \textbf{Industry Intern:} Product Maintenance Engineer Intern, Compound Co., Shenzhen (Summer 2023) -- sensor calibration \& IoT for medical devices
  \item \textbf{Languages:} English (fluent, in the US since 2016), Mandarin (native)
\end{itemize}

\end{document}
